\chapter{The contract}

\section{The target theorem}

\begin{theorem}[The exact Lean Pool target]\label{thm:mazur}
Let \(E/\Q\) be an elliptic curve.  The set underlying its rational torsion
subgroup has \texttt{ncard} at most \(16\):
\[
 \operatorname{ncard}\bigl(E(\Q)_{\mathrm{tors}}\bigr)\leq 16.
\]
Here \texttt{Set.ncard} is zero for an infinite set.  The project therefore
proves the cardinality consequence frozen in the Lean Pool challenge; it does
not silently claim that this proposition is logically equivalent to the
classical fifteen-group classification.
\uses{node:MT-FINAL-ASSEMBLY}
\leanpending
\end{theorem}

\begin{verbatim}
theorem torsion_ncard_le
    (E : WeierstrassCurve ℚ) [E.IsElliptic] :
    (AddCommGroup.torsion (E⁄ℚ).Point :
      Set (E⁄ℚ).Point).ncard ≤ 16
\end{verbatim}

\begin{remark}[Relation to the classical theorem]
Mazur's classical theorem says that \(E(\Q)_{\mathrm{tors}}\) is isomorphic
to \(\Z/n\Z\) for \(1\leq n\leq 10\) or \(n=12\), or to
\(\Z/2\Z\times\Z/2n\Z\) for \(1\leq n\leq 4\).  The point-order and
finite-abelian infrastructure on this roadmap follows that stronger
mathematics, but the percentage below is scoped to the exact Lean declaration
above.  Realizing every allowed group is outside this ledger.
\end{remark}

\begin{remark}[What the percentage means]\label{rem:percentage}
The canonical ledger is \texttt{coordination/program.json}.  Its 48 nodes
carry exactly 1000 points.  The imported, downstream-compiled baseline earns
50 points, so the honest theorem-completion estimate is
\[
  \frac{50}{1000}=5.0\%.
\]
The eight open contracts carry 113 points: 39 points are ordinary claimable
work and 74 points are nonexclusive research-open work.  An open statement
earns no credit.  Lines of Lean, repository size, and the number of
declarations are evidence of work performed, not a substitute for the
dependency-weighted percentage.
\end{remark}

\begin{definition}[Credit rules]\label{def:credit}
A proof node earns 70 percent when its isolated kernel proof compiles, 85
percent after API review, and 100 percent only after a real downstream
consumer compiles.  An infrastructure node earns 40 percent when its
definitions compile, 70 percent after mathematical sanity checks, and 100
percent after a downstream consumer compiles.  A theorem statement by itself
earns zero.
\end{definition}

\begin{center}
\begin{tabular}{lrr}
\hline
Stage & Weight & Earned now\\
\hline
Integrated baseline & 50 & 50\\
Finite-level endpoints & 100 & 0\\
Shared algebraic geometry and isogenies & 300 & 0\\
Prime-level infrastructure & 400 & 0\\
Mazur's prime-order argument & 100 & 0\\
Integration and hardening & 50 & 0\\
\hline
Total & 1000 & 50\\
\hline
\end{tabular}
\end{center}

\section{How to read the graph}

Every numbered item below corresponds one-for-one with a node in the JSON
ledger.  A \claimable{} node has an exact, compiling Lean statement whose only
hole is the registered \texttt{sorry}.  A \researchopen{} node has the same
exact contract, but its route is deliberately not prescribed and parallel
route-finding is welcome.  A \blocked{} node is useful planning information,
but should not be claimed until its missing nouns and immediate dependencies
exist.  Tau Ceti challenges live in a separately pinned package: they are
genuine upstream work, not imports into the current Mazur build.

\chapter{Integrated baseline}

\begin{milestone}[Imported sorry-free baseline:
\nodeid{MT-BASE-INTEGRATED} \weight{50}]\label{node:MT-BASE-INTEGRATED}
The pinned Clawristotle snapshot supplies the finite-group reductions,
low-torsion arithmetic, explicit descent infrastructure, and unconditional
exclusions of orders \(14,15,16,20,21,24,27\).  Its current public reduction
asks only for exclusions of the four composite orders \(18,25,35,49\) and
every prime order \(p\geq 11\).
\lean{MazurTorsion.rationalTorsion_orders_mem_cyclicOrders_of_remaining_obstructions}
\leanok
\end{milestone}

\chapter{Finite-level endpoints}

\section{\(X_1(11)\)}

\begin{theorem}[Five-coset bound:
\nodeid{MT-X11-COSET} \weight{12}, \claimable]\label{node:MT-X11-COSET}
Every rational point on the chosen model
\(y^2+y=x^3-x^2\) differs from one of the five visible torsion points by five
times a rational point.
\uses{node:MT-BASE-INTEGRATED}
\lean{MazurTheorem.Challenge.xOneEleven_fiveCosetBound}
\leanpending
\end{theorem}

\begin{verbatim}
theorem xOneEleven_fiveCosetBound :
    MazurTorsion.XOneEleven.FiveCosetBound := sorry
\end{verbatim}

Acceptance means replacing this one hole without weakening the statement,
then compiling the existing consumer
\texttt{five\_point\_classification\_of\_cosetBound}.  The intended route is
the explicit five-isogeny Kummer/Selmer calculation already prepared in
\texttt{XOneElevenDescent}.

\begin{milestone}[Order eleven join:
\nodeid{MT-X11-JOIN} \weight{2}, \blocked]\label{node:MT-X11-JOIN}
Feed the coset bound through the compiled \(X_1(11)\) reduction to exclude
exact rational order eleven.
\uses{node:MT-X11-COSET}
\leanpending
\end{milestone}

\section{\(X_1(13)\)}

\begin{theorem}[No noncuspidal point on \(X_1(13)\):
\nodeid{MT-X13-NONCUSP} \weight{26}, \researchopen]\label{node:MT-X13-NONCUSP}
The explicit order-thirteen genus-two model has no rational affine point
whose abscissa differs from both rational affine cusp abscissas \(0\) and
\(-1\).
\uses{node:MT-BASE-INTEGRATED}
\lean{MazurTheorem.Challenge.xOneThirteen_no_noncuspidal_point}
\leanpending
\end{theorem}

\begin{verbatim}
theorem xOneThirteen_no_noncuspidal_point
    (x y : ℚ) (hx0 : x ≠ 0) (hxnegOne : x ≠ -1)
    (hcurve : y ^ 2 =
      MazurTorsion.Kubert.orderThirteenHyperellipticPolynomial x) :
    False := sorry
\end{verbatim}

The repository already compiles the exact-order-\(13\) Tate-normal-form map
to this sextic.  Acceptance means compiling
\texttt{rationalPoint\_addOrderOf\_ne\_thirteen\_of\_noNoncuspidalPoint}.
The prepared Pell certificate and a genus-two Jacobian argument are one
route, but this research-open contract deliberately permits independent descents or
another kernel-checked rational-point classification.

\section{Order eighteen}

\begin{theorem}[No noncuspidal point on the sextic:
\nodeid{MT-X18-NONCUSP} \weight{18}, \researchopen]\label{node:MT-X18-NONCUSP}
If a rational pair \((x,y)\) lies on the compiled order-eighteen genus-two
model, then \(x=0\) or \(x=1\).
\uses{node:MT-BASE-INTEGRATED}
\lean{MazurTheorem.Challenge.xOneEighteen_no_noncuspidal_point}
\leanpending
\end{theorem}

\begin{verbatim}
theorem xOneEighteen_no_noncuspidal_point
    (x y : ℚ) (hx0 : x ≠ 0) (hx1 : x ≠ 1)
    (hcurve : y ^ 2 =
      MazurTorsion.Kubert.orderEighteenHyperellipticPolynomial x) :
    False := sorry
\end{verbatim}

The prepared route is the classical descent in the Eisenstein integers.
Acceptance includes the existing Tate-normal-form consumer, not merely an
isolated proof of an unrelated curve statement.

\section{The remaining composite orders}

\begin{theorem}[Exclude order twenty-five:
\nodeid{MT-O25-EXCLUDE} \weight{16}, \researchopen]\label{node:MT-O25-EXCLUDE}
No elliptic curve over \(\Q\) has a rational point of exact additive order
\(25\).
\uses{node:MT-BASE-INTEGRATED}
\lean{MazurTheorem.Challenge.no_rational_point_of_order_twentyFive}
\leanpending
\end{theorem}

\begin{verbatim}
theorem no_rational_point_of_order_twentyFive
    (E : WeierstrassCurve ℚ) [E.IsElliptic]
    (P : (E⁄ℚ).Point) :
    addOrderOf P ≠ 25 := sorry
\end{verbatim}

\begin{theorem}[Exclude order thirty-five:
\nodeid{MT-O35-EXCLUDE} \weight{14}, \researchopen]\label{node:MT-O35-EXCLUDE}
No elliptic curve over \(\Q\) has a rational point of exact additive order
\(35\).
\uses{node:MT-BASE-INTEGRATED}
\lean{MazurTheorem.Challenge.no_rational_point_of_order_thirtyFive}
\leanpending
\end{theorem}

\begin{verbatim}
theorem no_rational_point_of_order_thirtyFive
    (E : WeierstrassCurve ℚ) [E.IsElliptic]
    (P : (E⁄ℚ).Point) :
    addOrderOf P ≠ 35 := sorry
\end{verbatim}

\begin{theorem}[Exclude the order-forty-nine tower:
\nodeid{MT-O49-TOWER} \weight{10}, \claimable]\label{node:MT-O49-TOWER}
No elliptic curve over \(\Q\) has a rational point of exact additive order
\(49\).
\uses{node:MT-BASE-INTEGRATED}
\lean{MazurTheorem.Challenge.no_rational_point_of_order_fortyNine}
\leanpending
\end{theorem}

\begin{verbatim}
theorem no_rational_point_of_order_fortyNine
    (E : WeierstrassCurve ℚ) [E.IsElliptic]
    (P : (E⁄ℚ).Point) :
    addOrderOf P ≠ 49 := sorry
\end{verbatim}

The order-\(49\) route should expose the existing level-seven correspondence
API, construct the marked isogeny tower, and apply the compiled
\texttt{no\_noncuspidal\_correspondence\_point} theorem.

\begin{milestone}[Finite endpoint join:
\nodeid{MT-FINITE-JOIN} \weight{2}, \blocked]\label{node:MT-FINITE-JOIN}
Assemble the order \(11,13,18,25,35,49\) results behind one stable finite-level
interface.
\uses{node:MT-X11-JOIN,node:MT-X13-NONCUSP,node:MT-X18-NONCUSP,node:MT-O25-EXCLUDE,node:MT-O35-EXCLUDE,node:MT-O49-TOWER}
\leanpending
\end{milestone}

\chapter{Shared algebraic geometry and isogenies}

This 300-point stage is the main reusable layer.  Work should land in Mathlib
or Tau Ceti when its statement is not Mazur-specific.  The two compiled
upstream challenges below use Tau Ceti's own Lean and Mathlib pin; completing
them does not silently make that package compatible with the Mazur checkout.

\section{Divisors and cohomology}

\begin{theorem}[Finite support of orders:
\nodeid{MT-TC-A1-ORDER-SUPPORT} \weight{15}, \claimable]\label{node:MT-TC-A1-ORDER-SUPPORT}
A nonzero rational function on a Noetherian integral scheme has
nonzero order at only finitely many codimension-one points.
\uses{node:MT-BASE-INTEGRATED}
\lean{MazurTauCetiChallenge.finite_support_orderAt}
\leanpending
\end{theorem}

\begin{verbatim}
theorem finite_support_orderAt
    (X : Scheme.{u}) [IsIntegral X] [IsNoetherian X]
    (g : Additive X.functionFieldˣ) :
    (Function.support fun x : CodimensionOnePoint X =>
      TauCeti.AlgebraicGeometry.SchemeWeilDivisor.orderAt x g).Finite := sorry
\end{verbatim}

The same file contains \texttt{MazurTauCetiChallenge.orderSystem}, an
acceptance consumer which packages the theorem into Tau Ceti's
\texttt{OrderSystem}.

\begin{theorem}[Principal-divisor product formula:
\nodeid{MT-TC-A2-PRODUCT-FORMULA} \weight{15}, \blocked]\label{node:MT-TC-A2-PRODUCT-FORMULA}
Construct principal Weil divisors from the order system and prove the
degree-zero product formula on proper curves.
\uses{node:MT-TC-A1-ORDER-SUPPORT}
\leanpending
\end{theorem}

\begin{milestone}[Divisors to line bundles:
\nodeid{MT-TC-A3-DIVISOR-LINE-BUNDLE} \weight{18}, \blocked]\label{node:MT-TC-A3-DIVISOR-LINE-BUNDLE}
Relate Cartier and Weil divisors to invertible sheaves, including pullback,
degree, and linear equivalence on regular curves.
\uses{node:MT-TC-A2-PRODUCT-FORMULA}
\leanpending
\end{milestone}

\begin{milestone}[Coherent cohomology:
\nodeid{MT-TC-B1-COHERENT-COHOMOLOGY} \weight{35}, \blocked]\label{node:MT-TC-B1-COHERENT-COHOMOLOGY}
Define the finite-dimensional coherent cohomology groups needed for curves
and make their functoriality usable downstream.
\uses{node:MT-TC-A3-DIVISOR-LINE-BUNDLE}
\leanpending
\end{milestone}

\begin{theorem}[Riemann--Roch and Serre duality:
\nodeid{MT-TC-B2-RR-SERRE} \weight{25}, \blocked]\label{node:MT-TC-B2-RR-SERRE}
Prove the curve-level Riemann--Roch and Serre-duality statements required for
Picard and Abel--Jacobi constructions.
\uses{node:MT-TC-B1-COHERENT-COHOMOLOGY}
\leanpending
\end{theorem}

\begin{milestone}[Relative cohomology and base change:
\nodeid{MT-TC-C1-RELATIVE-COHOMOLOGY} \weight{30}, \blocked]\label{node:MT-TC-C1-RELATIVE-COHOMOLOGY}
Supply flat families of curves with cohomology-and-base-change results strong
enough for representability and Abel--Jacobi.
\uses{node:MT-TC-B1-COHERENT-COHOMOLOGY}
\leanpending
\end{milestone}

\begin{milestone}[Symmetric powers of curves:
\nodeid{MT-TC-C2-SYMMETRIC-POWERS} \weight{15}, \blocked]\label{node:MT-TC-C2-SYMMETRIC-POWERS}
Construct relative symmetric powers and their effective-divisor
interpretation.
\uses{node:MT-TC-A3-DIVISOR-LINE-BUNDLE,node:MT-TC-C1-RELATIVE-COHOMOLOGY}
\leanpending
\end{milestone}

\section{Picard schemes, Jacobians, and Abel--Jacobi}

\begin{milestone}[Picard functor:
\nodeid{MT-TC-D1-PICARD-FUNCTOR} \weight{35}, \blocked]\label{node:MT-TC-D1-PICARD-FUNCTOR}
Define the relative Picard functor with rigidification, sheafification, and
base-change behavior.
\uses{node:MT-TC-A3-DIVISOR-LINE-BUNDLE,node:MT-TC-C1-RELATIVE-COHOMOLOGY}
\leanpending
\end{milestone}

\begin{theorem}[Picard representability:
\nodeid{MT-TC-D2-PICARD-REPRESENTABILITY} \weight{45}, \blocked]\label{node:MT-TC-D2-PICARD-REPRESENTABILITY}
Represent the degree-zero Picard functor for a proper smooth geometrically
connected curve.
\uses{node:MT-TC-B2-RR-SERRE,node:MT-TC-C2-SYMMETRIC-POWERS,node:MT-TC-D1-PICARD-FUNCTOR}
\leanpending
\end{theorem}

\begin{theorem}[Dimension of a product of abelian varieties:
\nodeid{MT-TC-E0-PRODUCT-DIM} \weight{2}, \claimable]\label{node:MT-TC-E0-PRODUCT-DIM}
For abelian varieties \(A\) and \(B\) over a field,
\(\dim(A\times B)=\dim A+\dim B\).
\uses{node:MT-BASE-INTEGRATED}
\lean{MazurTauCetiChallenge.prod_dim}
\leanpending
\end{theorem}

\begin{verbatim}
theorem prod_dim {K : Type u} [Field K] (A B : AbelianVariety K) :
    (AbelianVariety.prod A B).dim = A.dim + B.dim := sorry
\end{verbatim}

\begin{milestone}[Jacobian variety and sanity checks:
\nodeid{MT-TC-E1-JACOBIAN-VARIETY} \weight{20}, \blocked]\label{node:MT-TC-E1-JACOBIAN-VARIETY}
Bundle \(\operatorname{Pic}^0\) as an abelian variety, prove that its dimension
is the genus, and recover a pointed elliptic curve from its Jacobian.
\uses{node:MT-TC-D2-PICARD-REPRESENTABILITY,node:MT-TC-E0-PRODUCT-DIM}
\leanpending
\end{milestone}

\begin{theorem}[Abel--Jacobi:
\nodeid{MT-TC-F1-ABEL-JACOBI} \weight{20}, \blocked]\label{node:MT-TC-F1-ABEL-JACOBI}
Construct the Abel--Jacobi map, prove its universal property and base-change
compatibility, and prove the needed closed-immersion result.
\uses{node:MT-TC-C1-RELATIVE-COHOMOLOGY,node:MT-TC-E1-JACOBIAN-VARIETY}
\leanpending
\end{theorem}

\section{Elliptic-curve isogeny infrastructure}

\begin{milestone}[Isogenies, quotients, duals, and Weil pairing:
\nodeid{MT-EC-ISOGENY-WEIL} \weight{25}, \planned]\label{node:MT-EC-ISOGENY-WEIL}
Build finite-subgroup quotients, dual isogenies, multiplication kernels, and
the Weil pairing, with explicit base-change tests.
\uses{node:MT-BASE-INTEGRATED}
\leanpending
\end{milestone}

\chapter{Prime-level infrastructure}

The 400 points in this chapter are why the project is not eight independent
``fill the hole'' problems.  They form parallel workstreams, but each
workstream needs reviewed nouns and consumer tests before the next layer is
claimable.

\section{N\'eron models}

\begin{milestone}[N\'eron models:
\nodeid{MT-NERON-BASE} \weight{40}, \blocked]\label{node:MT-NERON-BASE}
Define the N\'eron model and its mapping property in the generality needed by
Jacobians and isogenies.
\uses{node:MT-TC-E1-JACOBIAN-VARIETY,node:MT-EC-ISOGENY-WEIL}
\leanpending
\end{milestone}

\begin{milestone}[Component groups:
\nodeid{MT-NERON-COMPONENTS} \weight{30}, \blocked]\label{node:MT-NERON-COMPONENTS}
Develop identity components, component groups, and functoriality under
isogenies.
\uses{node:MT-NERON-BASE}
\leanpending
\end{milestone}

\begin{theorem}[Specialization:
\nodeid{MT-NERON-SPECIALIZATION} \weight{30}, \blocked]\label{node:MT-NERON-SPECIALIZATION}
Prove the torsion-specialization and component-group controls used in Mazur's
argument.
\uses{node:MT-NERON-COMPONENTS}
\leanpending
\end{theorem}

\section{Finite-flat group schemes}

\begin{milestone}[Finite-flat commutative group schemes:
\nodeid{MT-FFGS-BASIC} \weight{20}, \planned]\label{node:MT-FFGS-BASIC}
Establish the basic category, kernels, Cartier duality, and base change.
\uses{node:MT-BASE-INTEGRATED}
\leanpending
\end{milestone}

\begin{milestone}[Connected--\'etale sequence:
\nodeid{MT-FFGS-CONNECTED-ETALE} \weight{20}, \blocked]\label{node:MT-FFGS-CONNECTED-ETALE}
Construct connected and \'etale parts over the relevant local bases.
\uses{node:MT-FFGS-BASIC}
\leanpending
\end{milestone}

\begin{theorem}[Oort--Tate/Raynaud classification:
\nodeid{MT-FFGS-OORT-RAYNAUD} \weight{40}, \blocked]\label{node:MT-FFGS-OORT-RAYNAUD}
Classify the order-\(p\) finite-flat group schemes used to control reductions
of rational torsion.
\uses{node:MT-FFGS-CONNECTED-ETALE}
\leanpending
\end{theorem}

\section{Integral modular curves and Hecke algebra}

\begin{milestone}[The \(X_0(p)\) moduli problem:
\nodeid{MT-X0-MODULI} \weight{30}, \planned]\label{node:MT-X0-MODULI}
Formalize generalized elliptic curves with a cyclic subgroup and the
coarse/fine moduli interfaces needed later.
\uses{node:MT-BASE-INTEGRATED}
\leanpending
\end{milestone}

\begin{milestone}[Integral model:
\nodeid{MT-X0-INTEGRAL} \weight{30}, \blocked]\label{node:MT-X0-INTEGRAL}
Construct the integral model of \(X_0(p)\) and its special fibers.
\uses{node:MT-X0-MODULI}
\leanpending
\end{milestone}

\begin{theorem}[Cusps:
\nodeid{MT-X0-CUSPS} \weight{20}, \blocked]\label{node:MT-X0-CUSPS}
Classify rational cusps, their fields of definition, and specialization.
\uses{node:MT-X0-INTEGRAL}
\leanpending
\end{theorem}

\begin{milestone}[The modular Jacobian:
\nodeid{MT-X0-JACOBIAN} \weight{20}, \blocked]\label{node:MT-X0-JACOBIAN}
Instantiate the shared Jacobian and Abel--Jacobi APIs for \(X_0(p)\).
\uses{node:MT-X0-INTEGRAL,node:MT-TC-E1-JACOBIAN-VARIETY,node:MT-TC-F1-ABEL-JACOBI}
\leanpending
\end{milestone}

\begin{milestone}[Hecke correspondences:
\nodeid{MT-X0-HECKE} \weight{30}, \blocked]\label{node:MT-X0-HECKE}
Define the Hecke action on the modular Jacobian and prove its compatibility
with isogenies and reduction.
\uses{node:MT-X0-JACOBIAN,node:MT-EC-ISOGENY-WEIL}
\leanpending
\end{milestone}

\begin{theorem}[Eisenstein algebra:
\nodeid{MT-X0-EISENSTEIN-ALGEBRA} \weight{30}, \blocked]\label{node:MT-X0-EISENSTEIN-ALGEBRA}
Define the Eisenstein ideal and establish the algebraic relations used in the
quotient argument.
\uses{node:MT-X0-HECKE}
\leanpending
\end{theorem}

\begin{theorem}[Eisenstein quotient:
\nodeid{MT-X0-EISENSTEIN-QUOTIENT} \weight{40}, \blocked]\label{node:MT-X0-EISENSTEIN-QUOTIENT}
Construct the Eisenstein quotient and prove the required finiteness and
specialization properties.
\uses{node:MT-X0-CUSPS,node:MT-X0-EISENSTEIN-ALGEBRA,node:MT-NERON-SPECIALIZATION,node:MT-FFGS-OORT-RAYNAUD}
\leanpending
\end{theorem}

\section{Cyclotomic input}

\begin{theorem}[Cyclotomic unramifiedness:
\nodeid{MT-CYCLOTOMIC-UNRAMIFIED} \weight{20}, \planned]\label{node:MT-CYCLOTOMIC-UNRAMIFIED}
Package the ramification and class-group input for the relevant cyclotomic
fields in the form consumed by the prime argument.
\uses{node:MT-BASE-INTEGRATED}
\leanpending
\end{theorem}

\chapter{Mazur's prime-order argument}

\begin{theorem}[Semistability:
\nodeid{MT-PRIME-SEMISTABLE} \weight{10}, \blocked]\label{node:MT-PRIME-SEMISTABLE}
A rational point of prime order \(p\geq 17\) forces the required semistable
reduction behavior.
\uses{node:MT-NERON-SPECIALIZATION,node:MT-FFGS-OORT-RAYNAUD}
\leanpending
\end{theorem}

\begin{theorem}[Outside the identity component:
\nodeid{MT-PRIME-OUTSIDE-IDENTITY} \weight{10}, \blocked]\label{node:MT-PRIME-OUTSIDE-IDENTITY}
Control where the marked point specializes in the N\'eron model.
\uses{node:MT-PRIME-SEMISTABLE}
\leanpending
\end{theorem}

\begin{theorem}[Eisenstein specialization:
\nodeid{MT-PRIME-EISENSTEIN-SPECIALIZATION} \weight{20}, \blocked]\label{node:MT-PRIME-EISENSTEIN-SPECIALIZATION}
Map the marked modular point into the Eisenstein quotient and obtain the
specialization constraint.
\uses{node:MT-PRIME-OUTSIDE-IDENTITY,node:MT-X0-EISENSTEIN-QUOTIENT}
\leanpending
\end{theorem}

\begin{theorem}[Division field:
\nodeid{MT-PRIME-DIVISION-FIELD} \weight{15}, \blocked]\label{node:MT-PRIME-DIVISION-FIELD}
Derive the required structure and ramification of the prime-division field.
\uses{node:MT-PRIME-EISENSTEIN-SPECIALIZATION,node:MT-CYCLOTOMIC-UNRAMIFIED}
\leanpending
\end{theorem}

\begin{theorem}[Herbrand--Kummer step:
\nodeid{MT-PRIME-HERBRAND-KUMMER} \weight{10}, \blocked]\label{node:MT-PRIME-HERBRAND-KUMMER}
Apply the cyclotomic class-group input to rule out the forbidden extension
class.
\uses{node:MT-PRIME-DIVISION-FIELD}
\leanpending
\end{theorem}

\begin{theorem}[Split extension:
\nodeid{MT-PRIME-SPLIT-SEQUENCE} \weight{10}, \blocked]\label{node:MT-PRIME-SPLIT-SEQUENCE}
Split the resulting Galois module sequence and turn it into an isogeny
construction.
\uses{node:MT-PRIME-HERBRAND-KUMMER,node:MT-EC-ISOGENY-WEIL}
\leanpending
\end{theorem}

\begin{theorem}[Shafarevich finiteness input:
\nodeid{MT-PRIME-SHAFAREVICH} \weight{15}, \blocked]\label{node:MT-PRIME-SHAFAREVICH}
Supply the finiteness statement for the abelian varieties arising in the
isogeny iteration.
\uses{node:MT-TC-E1-JACOBIAN-VARIETY,node:MT-EC-ISOGENY-WEIL}
\leanpending
\end{theorem}

\begin{theorem}[Isogeny-chain contradiction:
\nodeid{MT-PRIME-ISOGENY-CHAIN} \weight{10}, \blocked]\label{node:MT-PRIME-ISOGENY-CHAIN}
Iterate the split isogeny construction and contradict Shafarevich finiteness,
excluding every rational point of prime order \(p\geq 17\).
\uses{node:MT-PRIME-SPLIT-SEQUENCE,node:MT-PRIME-SHAFAREVICH}
\leanpending
\end{theorem}

\chapter{Integration and hardening}

\begin{milestone}[Pin migration:
\nodeid{MT-PIN-MIGRATION} \weight{20}, \planned]\label{node:MT-PIN-MIGRATION}
Converge the Mazur, Mathlib, and upstream pins without weakening any theorem
or bypassing any checker.
\uses{node:MT-BASE-INTEGRATED}
\leanpending
\end{milestone}

\begin{milestone}[API integration:
\nodeid{MT-API-INTEGRATION} \weight{10}, \blocked]\label{node:MT-API-INTEGRATION}
Connect the finite endpoint and prime argument behind stable public APIs on
the converged pin.
\uses{node:MT-FINITE-JOIN,node:MT-PRIME-ISOGENY-CHAIN,node:MT-PIN-MIGRATION}
\leanpending
\end{milestone}

\begin{theorem}[Final assembly:
\nodeid{MT-FINAL-ASSEMBLY} \weight{15}, \blocked]\label{node:MT-FINAL-ASSEMBLY}
Apply the point-order theorem and the existing finite-abelian cardinality
reduction to obtain the exact Lean Pool declaration in
Theorem~\ref{thm:mazur}.
\uses{node:MT-API-INTEGRATION}
\leanpending
\end{theorem}

\begin{milestone}[Exposition and kernel audit:
\nodeid{MT-EXPOSITION-AUDIT} \weight{5}, \blocked]\label{node:MT-EXPOSITION-AUDIT}
Finish cross-linked documentation, compile all consumers, and audit every
MazurTorsion declaration against the axiom allowlist
\(\{\texttt{propext},\texttt{Quot.sound},\texttt{Classical.choice}\}\).
\uses{node:MT-FINAL-ASSEMBLY}
\leanpending
\end{milestone}

\chapter{Contribution protocol}

\section{Claiming work}

Only nodes marked \claimable{} or \researchopen{} have a stable compiled
statement today.  An ordinary claim should name the node ID, proposed route,
target repository, and expected first checkpoint.  Research intentions may
run in parallel until a route is selected; they do not grant exclusive
ownership of the problem.  For large tasks, split research, API design,
implementation, and integration into reviewable pull requests while keeping
the original acceptance consumer visible.

\section{Acceptance}

A contribution is accepted only when:
\begin{enumerate}
\item its exact registered statement remains unchanged, unless a roadmap
      review explicitly changes the contract before work begins;
\item the relevant Lean package builds on its pinned toolchain;
\item no \texttt{sorry}, \texttt{admit}, \texttt{native\_decide}, custom
      axiom, \texttt{unsafe}, or \texttt{partial} enters completed mathematics;
\item the standard axiom audit reports only \texttt{propext},
      \texttt{Quot.sound}, and \texttt{Classical.choice};
\item the registered downstream consumer compiles; and
\item \texttt{coordination/program.json}, the blueprint, and the public
      progress display are updated together.
\end{enumerate}

\section{Why this is crowdsourceable}

The graph does not claim that all remaining work can be split into independent
one-week theorem holes.  It makes a more useful claim: the program can be
crowdsourced as a sequence of independently reviewed interfaces.  At launch,
39 of 1000 points are ordinary claimable work and 74 are nonexclusive
research-open contracts.  The other nodes become safely parallel only as
divisor, cohomology, Jacobian, isogeny, N\'eron-model, finite-flat, and
modular-curve APIs acquire compiled acceptance consumers.  That is a staged
crowdsourcing program, not a promise that 95 percent of the theorem is already
shovel-ready.
